\documentclass[10pt]{amsart}
\usepackage[a4paper,margin=22mm]{geometry}
\usepackage[no-math]{fontspec}
\setmainfont{Latin Modern Roman}[SmallCapsFont={lmromancaps10-regular.otf}]
\usepackage{amsmath,amssymb,amsthm,mathtools,hyperref,graphicx,xfrac,etoolbox}
\hypersetup{hidelinks}
\usepackage[scr=rsfs,cal=euler]{mathalfa}
\newcommand{\ie}{i.e.}
\newcommand{\eg}{e.g.}
\newcommand{\cf}{cf.}
\newcommand{\resp}{resp.}
\newcommand{\skpt}{\par\smallskip\noindent}
\let\Subsection\subsection
\let\Subsubsection\subsubsection
\emergencystretch=3em
\newenvironment{enumeratea}
{\bgroup\def\theenumi{\alph{enumi}}\begin{enumerate}}
{\end{enumerate}\egroup}

\newenvironment{smallpmatrix}{\big(\begin{smallmatrix}}{\end{smallmatrix}\big)}

\newenvironment{Smallpmatrix}{\Big(\begin{smallmatrix}}{\end{smallmatrix}\Big)}

\newcommand\mto{\mathchoice{\longmapsto}{\mapsto}{\mapsto}{\mapsto}}
\newcommand\hto{\mathrel{\lhook\joinrel\to}}
\newcommand{\onto}{\mathchoice{\to\hspace*{-5mm}\to}{\mathrel{\hbox{$\to\hspace*{-2.85mm}\to$}}}{\to\hspace*{-3mm}\to}{\to\hspace*{-3mm}\to}}
\newcommand\from{\mathchoice{\longleftarrow}{\leftarrow}{\leftarrow}{\leftarrow}}
\def\hfrom{\mathrel{\from\joinrel\rhook}}
\newcommand{\onfrom}{\mathchoice{\from\hspace*{-5mm}\from}{\from\hspace*{-2.85mm}\from}{\from\hspace*{-3mm}\from}{\from\hspace*{-3mm}\from}}
\newcommand{\To}[1]{\mathchoice{\xrightarrow{\textstyle\kern4pt#1\kern3pt}}{\stackrel{#1}{\longrightarrow}}{}{}}
\newcommand{\quand}{\quad\text{and}\quad}
\newcommand{\sbullet}{{\scriptscriptstyle\bullet}}

\newcommand{\Psfrac}[2]{(\sfrac{#1}{#2})}
\newcommand{\psfrac}[2]{\sfrac{(#1)}{#2}}
\newcommand{\spfrac}[2]{\sfrac{#1}{(#2)}}
\newcommand{\pspfrac}[2]{\sfrac{(#1)}{(#2)}}

\let\backslash\setminus
\let\wt\widetilde
\let\ovl\overline
\let\un\underline

\usepackage[all,cmtip]{xy}\let\labelstyle\textstyle
\newdir^{ (}{{}*!/-5pt/\dir^{(}}
\newdir_{ (}{{}*!/-5pt/\dir_{(}}

\renewcommand{\theenumi}{\roman{enumi}}

\newcommand{\defeq}{\stackrel{\textrm{\tiny{\upshape{def}}}}{=}}

\newcommand{\F}{\mathbb{F}}
\newcommand{\Q}{\mathbb{Q}}
\newcommand{\Z}{{\mathbb Z}}

\newcommand{\cI}{\mathcal{I}}
\newcommand{\cJ}{{\mathcal{J}}}
\newcommand{\cO}{\mathcal{O}}
\newcommand{\cS}{\mathcal{S}}
\newcommand{\cW}{\mathcal{W}}
\newcommand{\cZ}{\mathcal{Z}}

\newcommand{\D}{\mathscr D}
\renewcommand{\P}{\mathscr P}

\newcommand{\EE}{\mathrm{E}}
\newcommand{\sgn}{\mathrm{sgn}}
\renewcommand{\ss}{\mathrm{ss}}

\newcommand{\fa}{\mathfrak{a}}
\newcommand{\fb}{\mathfrak{b}}
\newcommand{\fm}{\mathfrak{m}}\let\m\fm
%\newcommand{\n}{\mathfrak{n}}

\newcommand{\Fp}{\mathbb{F}_{p}}
\newcommand{\Fq}{\mathbb{F}_{q}}
\newcommand{\Qp}{\mathbb{Q}_{p}}
\newcommand{\rbar}{\overline{r}}
\newcommand{\rhobar}{\overline{\rho}}
\newcommand{\brho}{\overline{\rho}}

\newcommand{\s}{^\times}
\newcommand{\ppar}[1]{(\mkern-3mu(#1)\mkern-3mu)}
\newcommand{\bbra}[1]{[\mkern-3mu[#1]\mkern-3mu]}

\DeclareMathOperator{\Ann}{Ann}
\DeclareMathOperator{\coker}{coker}
\DeclareMathOperator{\cosoc}{cosoc}
\DeclareMathOperator{\Ext}{Ext}
\DeclareMathOperator{\Gal}{Gal}
\DeclareMathOperator{\GL}{GL}
\DeclareMathOperator{\gr}{gr}
\DeclareMathOperator{\Hom}{Hom}
\DeclareMathOperator{\HOM}{HOM}
\DeclareMathOperator{\im}{im}
\DeclareMathOperator{\Ind}{Ind}
\DeclareMathOperator{\Inj}{Inj}
\DeclareMathOperator{\JH}{\mathrm{JH}}
\DeclareMathOperator{\PGL}{PGL}
\DeclareMathOperator{\Proj}{Proj}
\DeclareMathOperator{\rad}{rad}
\DeclareMathOperator{\soc}{soc}
\DeclareMathOperator{\Tor}{Tor}

\let\ve\varepsilon
\let\vp\varphi
\AtBeginDocument{\let\ra\to}

\newcommand{\smatr}[4]{\bigl(\begin{smallmatrix} {#1}& {#2}\\ {#3}&{#4}\end{smallmatrix}\bigr)}
\newcommand{\xto}[1][]{\To{#1}}
\newcommand{\simto}{\xrightarrow{\,\sim\,}}\let\congto\simto

\theoremstyle{plain}
\newtheorem{prop}{Proposition}[subsection]
\newtheorem{lem}[prop]{Lemma}
\newtheorem{thm}[prop]{Theorem}
\newtheorem{cor}[prop]{Corollary}
\newtheorem{prop1}{Proposition}[section]
\newtheorem{lem1}[prop1]{Lemma}
\newtheorem{thm1}[prop1]{Theorem}
\newtheorem{cor1}[prop1]{Corollary}

\theoremstyle{definition}
\newtheorem{rem}[prop]{Remark}
\newtheorem{rem1}[prop1]{Remark}

\newcommand{\loccit}{the cited passage}
\AtBeginDocument{
\let\sourceLocalRef\ref
\renewcommand{\ref}[1]{\ifstrequal{#1}{sec:results}{\href{https://doi.org/10.5802/jep.346}{1.1 (source)}}{\ifstrequal{#1}{sec:subquot-ss}{\href{https://doi.org/10.5802/jep.346}{5 (source)}}{\ifstrequal{#1}{eq:1st-term}{\href{https://doi.org/10.5802/jep.346}{6 (source)}}{\ifstrequal{#1}{cor:gr-pi'}{\href{https://doi.org/10.5802/jep.346}{6.1.7 (source)}}{\sourceLocalRef{#1}}}}}}
}
\title[Graded quotients in the nonsplit case]{The associated graded module of a quotient of a generic nonsplit mod p GL2 representation}
\author{Christophe Breuil; Florian Herzig; Yongquan Hu; Stefano Morra; Benjamin Schraen}
\date{Published version; DOI 10.5802/jep.346}
\begin{document}
\maketitle
\noindent\textit{Source notation (Section 1.1).} $K$ is the unramified finite extension of $\Qp$ considered in the source, $f=[K:\Qp]$, $q=p^f$, and $\F$ is a sufficiently large finite extension of $\Fp$. The characters $\omega_f$ and $\omega_{2f}$ are Serre's fundamental inertia characters of levels $f$ and $2f$, and $\omega$ is the mod $p$ cyclotomic character. This item retains the abstract assumptions of Section 3.1 and counts only Theorem 6.1.1.
\section*{Local notation and assumptions}
\setcounter{section}{1}\setcounter{subsection}{3}\setcounter{equation}{9}
\subsection{Notation and preliminaries}
\label{sec:notation}
\label{sec:preliminaries}

We normalize local class field theory so that uniformizers correspond to geometric Frobenius elements. We~fix an embedding \hbox{$\kappa_0:\F_q\hto \F$} and set $\kappa_j\defeq \kappa_0\circ\varphi^j$, where $\varphi$ is the arithmetic Frobenius on $\F_q$. Given $J\subset\{0,\dots,f-1\}$ we define $J^c\defeq \{0,1,\dots,f-1\} \setminus J$. We~let
\[
I\defeq \begin{pmatrix}\cO_K^\times&\cO_K\\p\cO_K&\cO_K^\times\end{pmatrix} \subset \GL_2(\cO_K)
\]
denote the (upper) Iwahori subgroup of $\GL_2(K)$, $I_1$ the pro-$p$ radical of $I$, $Z_1$ the center of $I_1$, and $K_1\defeq 1+p{\rm M}_2(\cO_K)\subset I_1$. We~let $\Gamma\defeq \GL_2(\F_q)\cong \GL_2(\cO_K)/K_1$.

Let $\brho:\Gal(\ovl K/K)\to \GL_2(\F)$ be a continuous representation. We~will say that $\rhobar$ is \emph{$n$-generic} for some integer $n \ge 0$ if, up to twist, $\rhobar|_{I_K}^{\ss}\not\cong \omega^{\pm 1}\oplus 1$ and either (using the notation of Section~\ref{sec:results})\vspace*{-3pt}
\begin{multline}\label{eq:rhobar-generic-red}
\rhobar|_{I_K} \cong \begin{pmatrix}\omega_f^{\sum_{j=0}^{f-1} (r_j+1) p^j} & * \\ & 1\end{pmatrix}\\
\text{with $n \leq r_j \leq p-3-n$ for all $j$ such that $0 \le j \le f-1$,}
\end{multline}
or\vspace*{-3pt}
\begin{multline}\label{eq:rhobar-generic-irred}
\rhobar|_{I_K} \cong \begin{pmatrix}\omega_{2f}^{\sum_{j=0}^{f-1} (r_j+1) p^j} & \\ & \omega_{2f}^{p^f(\text{same})}\end{pmatrix} \\ \text{with\ }
\begin{cases}
n \leq r_j \leq p-3-n &\text{for $0 < j \le f-1$,} \\ n+1 \leq r_0 \leq p-2-n & \text{for $j = 0$.}
\end{cases}
\end{multline}

In particular, if~$\rhobar$ is $n$-generic then it is $n$-generic in the sense of \cite[Def.\,2.3.4]{BHHMS1}, and $\rhobar$ is $0$-generic precisely when $\rhobar$ is generic in the sense of \cite[Def.\,11.7]{BP} (note that the condition $\rhobar|_{I_K}^{\ss}\not\cong \omega^{\pm 1}\oplus 1$ precisely rules out that\vspace*{-3pt}
\[
(r_0,\dots,r_{f-1})\in\{(0,\dots,0),(p-3,\dots,p-3)\}
\]
when $\rhobar$ is reducible).

Attached to a $0$-generic $\rhobar$ we have a set $W(\rhobar)$ of Serre weights, \ie irreducible representations of $\Gamma$ over $\F$, defined in \cite[\S 3]{BDJ}, and a multiplicity-free, finite-length $\Gamma$-representation $D_0(\rhobar)$ over $\F$, defined in \cite[\S 13]{BP}, which is of the form $D_0(\rhobar)=\bigoplus_{\tau\in W(\rhobar)}D_{0,\tau}(\rhobar)$, where each $D_{0,\tau}(\rhobar)$ is indecomposable (and multiplicity-free) with socle the Serre weight $\tau$ (\cite[\S 15]{BP}).

Suppose that $\rhobar$ is $0$-generic. Recall the set $\P$ parametrizing $\JH(D_0(\rhobar)^{I_1})$, see \cite[\S 4]{breuil-buzzati} (and denoted there by $\mathscr{PD}$, \resp $\mathscr{PID}$, if~$\rhobar$ is reducible, \resp irre\-ducible). Recall also the subset $\D \subset \P$ parametrizing (the $I_1$-invariants of) the set of Serre weights in $W(\brho)$ (denoted in \loccit\ by $\D$ or $\mathscr{ID}$ if $\rhobar$ is reducible or irreducible, respectively).

We let $\D^\ss\subset \P^\ss$ denote the corresponding sets for the semisimplification $\rhobar^\ss$ of~$\rhobar$, so~$\P \subset \P^\ss$ and $\D \subset \D^\ss$.

Since we will be using this many times, we~recall more precisely that, if $\brho$ is reducible, $\P^\ss$ denotes the set of $f$-tuples $(\lambda_0(x_0),\dots,\lambda_{f-1}(x_{f-1}))$ such that:
\begin{enumerate}
\item $\lambda_j(x_j) \in \{x_j,x_j+1,x_j+2,p-3-x_j,p-2-x_j,p-1-x_j\}$;
\item if $\lambda_j(x_j) \in \{x_j,x_j+1,x_j+2\}$, then $\lambda_{j+1}(x_{j+1}) \in \{x_{j+1},x_{j+1}+2,p-2-x_{j+1}\}$;
\item if $\lambda_j(x_j) \in \{p-3-x_j,p-2-x_j,p-1-x_j\}$, then
\[
\lambda_{j+1}(x_{j+1}) \in \{x_{j+1}+1,p-3-x_{j+1},p-1-x_{j+1}\}
\]
\end{enumerate}
and $\D^\ss$ is the subset such that $\lambda_j(x_j) \in \{x_j,x_j+1,p-3-x_j,p-2-x_j\}$.
Moreover, there exists a unique subset $J_{\rhobar} \subset \{0,\dots,f-1\}$ such that\vspace*{-3pt}
\begin{align}
\D &= \bigl\{\lambda \in \D^\ss: \lambda_j(x_j) \in \{x_j+1,p-3-x_j\} \Rightarrow j \in J_{\rhobar}\bigr\},\notag \\
\P &= \bigl\{\lambda \in \P^\ss: \lambda_j(x_j) \in \{x_j+2,p-3-x_j\} \Rightarrow j \in J_{\rhobar}\bigr\}.\label{eq:P}
\end{align}
In particular, $|W(\brho)| = 2^{|J_{\brho}|}$.

For $\lambda\in\mathscr{P}$ we denote by $\chi_\lambda$ the character of $H$ corresponding to $\lambda$.
(More precisely, in \cite[\S 4]{breuil-buzzati} a Serre weight $\sigma_\lambda$ is associated to $\lambda \in \P$ and $\chi_\lambda$ is the action of $H=I/I_1$ on the $1$-dimensional subspace $\sigma_\lambda^{I_1}$.)

Set
\begin{equation}\label{eq:J-lambda}
J_{\lambda}\defeq
\bigl\{j\in\{0,\dots,f-1\}: \lambda_j(x_j)\in \{x_j+1,x_j+2,p-3-x_j\}\bigr\}
\end{equation}
and let $\ell(\lambda) \defeq |J_\lambda|.$
By \cite[\S 11]{BP}, the map $\lambda \mto J_\lambda$ induces a bijection between $\mathscr{D}^\ss$ and the set of subsets of $\{0,\dots,f-1\}$.
Sometimes we will abuse notation and write $J_\tau \defeq J_\lambda$ and $\ell(\tau) \defeq \ell(\lambda)$ if $\tau \in W(\brho^\ss)$ is parametrized by $\lambda \in \D^\ss$.
Given $\lambda\in \D^\ss$ with corresponding subset $J=J_\lambda\subset\{0,\dots,f-1\}$ we write $\delta(\lambda)\in \D^\ss$ for the $f$-tuple defined by $\delta(\lambda)_j\defeq \lambda_{j+1}$ for all $j\in\{0,\dots,f-1\}$, and $\delta(J)\subset\{0,\dots,f-1\}$ for the subset corresponding to $\delta(\lambda)$.

As in \cite[\S 1]{BP}, given $f$ integers $r_0,\dots,r_{f-1}\in\{0,\dots,p-1\}$, we denote by $(r_0,\dots,r_{f-1})$ the Serre weight
\begin{equation*}
\mathrm{Sym}^{r_0}\F^2\otimes_{\F}(\mathrm{Sym}^{r_1}\F^2)^{\rm Fr}\otimes\cdots\otimes_{\F}(\mathrm{Sym}^{r_{f-1}}\F^2)^{{\rm Fr}^{f-1}},
\end{equation*}
where $\GL_2(\Fq)$ acts on $(\mathrm{Sym}^{r_j}\F^2)^{{\rm Fr}^j}$ via $\kappa_j:\Fq\hto \F$.
Following \cite[\S 2]{HuWang2}, we~say that a Serre weight is \emph{$m$-generic} for some integer $m\geq 0$ if $\sigma\cong (r_0,\dots,r_{f-1})$ up to twist, where $m\leq r_j\leq p-2-m$ for all $j\in \{0,\dots,f-1\}$. We~say that an $\F$-valued character $\chi$ of $I$ is \emph{$m$-generic} if $\chi=\sigma^{I_1}$ for some $m$-generic Serre weight $\sigma$.

For any smooth character $\chi: I \to \F\s$, we define $\chi^s\defeq \chi(\Pi (\cdot) \Pi^{-1})$ with $\Pi\defeq \begin{smallpmatrix}0&1\\ p&0\end{smallpmatrix}$.
If $\sigma$ is a Serre weight, we~write $\chi_\sigma$ for the character of $I/I_1$ on $\sigma^{I_1}$ and $\sigma^{[s]}$ for the unique Serre weight distinct from $\sigma$ such that $\chi_{\sigma^{[s]}}=\chi_{\sigma}^s$.
Finally, if~$\chi,\chi': I\to\F^\times$ are smooth characters such that $\Ext^1_{I/Z_1}(\chi',\chi)\neq 0$ we let $E_{\chi,\chi'}$ denote the unique nonsplit extension of $\chi'$ by $\chi$, \ie $0\to\chi\to E_{\chi,\chi'}\to\chi'\to 0$.
(The uniqueness follows from \cite[Lem.\,2.4]{yongquan-algebra}.)

Let $\fm_{K_1}$ denote the maximal ideal of {the Iwasawa algebra} $\F\bbra{K_1/Z_1}$ and let
\[
\wt{\Gamma}\defeq \F\bbra{\GL_2(\cO_K)/Z_1}/\fm_{K_1}^2
\]
(a finite-dimensional $\F$-algebra). We~use the terms ``$\wt{\Gamma}$-representations'' and ``$\wt{\Gamma}$\nobreakdash-mod\-ules'' interchangeably.

We write $\wt{D}_{0}(\brho)$ for the finite-dimensional $\wt{\Gamma}$-representation over $\F$ constructed in \cite[Prop.\,4.3]{HuWang2}. It~is the unique (up to isomorphism) $\wt{\Gamma}$-representation that is maximal with respect to the following two properties:
\begin{enumerate}
\item $\soc_{\wt{\Gamma}}\wt{D}_0(\brho)\cong\bigoplus_{\sigma\in W(\brho)}\sigma$;
\item any Serre weight of $W(\brho)$ occurs in $\wt{D}_0(\brho)$ with multiplicity one.
\end{enumerate}

Let $\Lambda \defeq \F\bbra{I_1/Z_1}$, a complete Noetherian local ring with maximal ideal $\m \defeq \m_{I_1/Z_1}$, and let $\gr(\Lambda) \defeq \gr_\m(\Lambda) = \bigoplus_{n\geq 0}\m^n/\m^{n+1}$ be the associated graded ring of $\Lambda$ with respect to the $\fm$-adic filtration on $\Lambda$. The rings $\Lambda$ and $\gr(\Lambda)$ are Auslander regular (see \cite[Th.\,5.3.4]{BHHMS1} with \cite[Th.\,III.2.2.5]{LiOy}). Recall (\cite[\S 3.1]{BHHMS2}) that we have an isomorphism of (noncommutative) algebras
\begin{equation}\label{eq:grlambda}
\gr(\Lambda)\cong \bigotimes_{j\in\{0,\dots,f-1\}}\F\langle y_j,z_j,h_j\rangle
\end{equation}
with relations $[y_j,z_j]=h_j$, $[h_j,z_i]=[y_i,h_j]=0$ for all $i,j\in\{0,\dots,f-1\}$. We~use increasing filtrations throughout, \ie $F_n \Lambda = \m^{-n}$ for $n \le 0$, and the degrees of~$y_j$ and~$z_j$ (\resp $h_j$) are $-1$ (\resp $-2$).
Define the graded ideal \hbox{$J \defeq (h_j,y_jz_j: 0 \!\le\! j \!\le\! f-1)$} of~$\gr(\Lambda)$. As~in \cite[\S 3]{BHHMS2}, we define
\[
R \defeq \gr(\Lambda)/(h_j: 0 \le j \le f-1) \cong \F[y_j,z_j: 0 \le j \le f-1],
\]
which is the largest commutative quotient of $\gr(\Lambda)$. We~also define the following quotient of $R$:
\[
\ovl R \defeq \gr(\Lambda)/J \cong R/(y_jz_j: 0 \le j \le f-1).
\]

We recall from \cite[Def.\,3.57]{BHHMS2} that, given $\lambda\in\mathscr{P}$, we have an associated ideal $\mathfrak{a}(\lambda)=(t_0,\dots,t_{f-1})$ of $R$, where the $t_j = t_j(\lambda)$ are defined as follows:
\begin{equation}
\label{eq:id:al}
t_j\defeq \begin{cases}
z_j& \text{if } \lambda_j(x_j)\in\{x_j,p-3-x_j\} \ \text{and}\ j\in J_{\brho}\\
y_j& \text{if } \lambda_j(x_j)\in\{x_j+2,p-1-x_j\} \ \text{and}\ j\in J_{\brho}\\
y_jz_j & \text{if } \lambda_j(x_j)\in \{x_j,p-1-x_j\} \ \text{and}\ j\notin J_{\brho}\\
y_jz_j& \text{if } \lambda_j(x_j)\in \{x_j+1,p-2-x_j\}.\end{cases}
\end{equation}
Note that $(y_jz_j: 0 \le j \le f-1) \subset \fa(\lambda)$, so~we often think of $\fa(\lambda)$ as an ideal of $\ovl R$.

Let
\[
H\defeq \begin{pmatrix}\F_q^\times&0\\0&\F_q^\times\end{pmatrix}\cong I/I_1.
\]
We~write $\alpha_j:H\to\F^\times$ for the character defined by $\begin{smallpmatrix}a&0\\0&d\end{smallpmatrix}\mto \kappa_j(ad^{-1})$. We~recall that, for any $j\in\{0,\dots,f-1\}$, the element $y_j$ (\resp $z_j$, \resp $h_j$) is an $H$-eigenvector with associated eigencharacter $\alpha_j$ (\resp $\alpha_j^{-1}$, \resp the trivial character).

Suppose that $H'$ is a compact $p$-adic analytic group and that $\pi_1$, $\pi_2$ are smooth representations of $H'$ over $\F$. We~write $\Ext^i_{H'}(\pi_1,\pi_2)$ for the $i$-th Ext group computed in the category of smooth representations of $H'$ over $\F$.
Dually, the functors $\Tor_i^{\F\bbra{H'}}(\pi_1^\vee,\pi_2^\vee)$ and $\Ext^i_{\F\bbra{H'}}(\pi_1^\vee,\pi_2^\vee)$ are computed in the abelian category of pseudocompact $\F\bbra{H'}$-modules.
(See, for example, \cite[\S 2]{emerton-ordII}.)

If $\sigma$ has finite length, we~write $\JH(\sigma)$ for its set of irreducible constituents up to isomorphism.

Throughout this paper, if~$R$ is a filtered (\resp graded) ring, a morphism of filtered (\resp graded) $R$-modules $f:M\to N$ will always be a \emph{filtered (\resp graded) morphism of degree zero}, \ie satisfying $f(M_i)\subset N_i$ for all $i\in \Z$.
If $R$ is any ring and $M$ any left $R$-module, we~recall that $\Ext^{i}_R(M,R)$, for $i\in \Z_{\ge 0}$, is a right $R$-module (for $i=0$ the right $R$-action is given by $(f r)(m)\defeq f(m)r$ for $r\in R$, $f\in \Hom_R(M,R)$ and $m\in M$), and we use the notation $\EE^i_R(M)\defeq \Ext^{i}_R(M,R)$.

\setcounter{section}{2}
\section{Abstract setting}
\label{sec:abstract-setting}

Let $\rhobar:\Gal(\ovl K/K)\ra\GL_2(\F)$ be a continuous $0$-generic representation as in Section~\ref{sec:notation} and let $\pi$ denote a smooth representation of $\GL_2(K)$ over $\F$. In~this section we introduce and study certain assumptions on $\pi$ (relative to $\brho$) that play a key role in our work.

\subsection{Assumptions}
\label{subsec:assumptions}

\emph{From now until the end of this paper}, we~assume that $\pi$ satisfies assumptions (i) (with $r = 1$) and (ii) in \cite[\S 3.3.2]{BHHMS2} and assumption (iv) (with $r = 1$) in \cite[\S 2.1]{BHHMS4}, \ie

\begin{enumerate}
\item\label{it:assum-i} we have $\pi^{K_1}\cong D_0(\rhobar)$ as $\GL_2(\cO_K)$-representations (in~particular, $\pi$ is admissible) and $\pi$ has central character $\det(\rhobar)\omega^{-1}$;
\item\label{it:assum-ii} for any $\lambda\in\P$ we have $[\pi[\fm^3]:\chi_\lambda]=[\pi[\fm]:\chi_\lambda]$;

\setcounter{enumi}{3}
\item\label{it:assum-iv} for any smooth character $\chi:I\to \F^{\times}$ and any $i \ge 0$, $\Ext^i_{I/Z_1}(\chi,\pi)\neq0 $ only if
$[\pi[\m]:\chi]\neq0$, in which case
\[
\dim_{\F}\Ext^i_{I/Z_1}(\chi,\pi)=\binom{2f}{i}.
\]
\end{enumerate}

For later reference we also recall assumption (iii) of \cite[\S 3.3.5]{BHHMS2}, \emph{though we will not assume it until Section~\ref{sec:subquot-ss}}:
\begin{enumerate}\setcounter{enumi}{2}
\item\label{it:assum-iii} there is a $\GL_2(K)$-equivariant isomorphism of $\Lambda$-modules
\[
\EE^{2f}_\Lambda(\pi^\vee) \cong \pi^\vee\otimes (\det(\rhobar)\omega^{-1}),
\]
where $\EE^{2f}_\Lambda(\pi^\vee)$ is endowed with the $\GL_2(K)$-action defined in \cite[Prop.\,3.2]{Ko}.
\end{enumerate}

Finally, we~introduce a further assumption which will be used only in Section~\ref{sec:subquot-non-ss} (namely to verify equation~\eqref{eq:1st-term} in the introduction).
\begin{enumerate}{\setcounter{enumi}{4}}
\item
\label{it:assum-v}

We have
\[
\dim_\F\Tor_1^{\gr(\Lambda)}(\gr(\Lambda)/\ovl\m^3,\gr_\m(\pi^{\vee})) = \dim_\F \Tor_1^{\Lambda}(\Lambda/\m^3,\pi^{\vee}),
\]
where $\ovl\m = (y_j,z_j: 0 \le j \le f-1)$ denotes the unique maximal graded ideal of $\gr(\Lambda)$ (see (\ref{eq:grlambda})).
\end{enumerate}
\setcounter{section}{5}
\section{On the structure of subquotients of \texorpdfstring{$\pi$}{pi} in the non-semisimple case}
\label{sec:subquot-non-ss}

We prove many results on the structure of subquotients of $\pi$ as $I$- and $\GL_2(\cO_K)$-representations.

From now on $\pi$ denotes an admissible smooth representation of $\GL_2(K)$ over $\F$ satisfying assumptions \eqref{it:assum-i}--\eqref{it:assum-v} of Section~\ref{sec:abstract-setting}, with underlying Galois representation $\rhobar$ which is nonsplit reducible and {0}-generic.

The main results of this section include the description of the $I_1$- and $K_1$-invariants as well as of the $\GL_2(\cO_K)$-socle of any subquotient of $\pi$.
These results all depend on determining the $I_1$-socle filtration of any subquotient $\pi'$ of $\pi$ (equivalently, the associated graded module of $\pi'^\vee$ for the $\m$-adic filtration), which is the subject of Section~\ref{sec:grad-module-subq}.

We again suppose that $\pi_1 \subset \pi$ is a subrepresentation of $\pi$ and let $\pi_2\defeq \pi/\pi_1$.
Let $i_0 \defeq i_0(\pi_1) \in \{-1,\dots,f\}$, \cf \cite[Th.\,4.3.15]{BHHMS4}.
To simplify notation, for $\lambda \in \P$ we let $d_\lambda \defeq \max\{i_0+1-|J_\lambda|,0\}$.

\subsection{The graded module of subquotient representations of \texorpdfstring{$\pi$}{pi}}
\label{sec:grad-module-subq}

We describe $\gr_\m(\pi'^\vee)$, where $\pi'$ is any subquotient of $\pi$ (Corollary \ref{cor:gr-pi'}). We~start with quotients $\pi_2 = \pi/\pi_1$ of $\pi$:

\begin{thm}\label{thm:gr-pi2}
Assume that $\brho$ is $\max\{9,2f+3\}$-generic. We have an isomorphism of graded $\gr(\Lambda)$-modules with compatible $H$-actions,
\begin{equation}\label{eq:gr-pi2}
\gr_\m(\pi_2^\vee) \cong \bigoplus_{\lambda \in \P} \chi_\lambda^{-1} \otimes \frac{\mathfrak{a}_1^{i_0}(\lambda)}{\mathfrak{a}(\lambda)}(-d_\lambda).
\end{equation}
\end{thm}

We recall the ideal $\fa_1^{i_0}(\lambda)$ of $\ovl R$ from~\cite[(77)]{BHHMS4}.
Let
\[
J_1 \defeq \{ j \in J_{\rhobar}^c: \lambda_j(x_j) = p-1-x_j \}\quand J_2 \defeq \{ j \in J_{\rhobar}^c: \lambda_j(x_j) = x_j \}.
\]
If $i_0 \ge |J_\lambda|$, then $\fa_1^{i_0}(\lambda)$ is the ideal generated by $\fa(\lambda)$ (\cf \eqref{eq:id:al}) and all
\[
\prod_{j \in J_1'} y_j \prod_{j \in J_2'} z_j, \quad\text{where }J_1' \subset J_1,\; J_2' \subset J_2,\;|J'_1|+|J'_2| = i_0+1-|J_\lambda|.
\]
Otherwise, $\fa_1^{i_0}(\lambda) = \ovl R$. In~particular, $\mathfrak a_1^{i_0}(\lambda) = \mathfrak a(\lambda)$ if $|J_1|+|J_2| < i_0+1-|J_\lambda|$.

The grading shifts in~\eqref{eq:gr-pi2} are such that all nonzero direct summands contribute in degree 0, but vanish in degree 1.
Note also from the definitions that $\mathfrak{a}_1^{i_0}(\lambda)/\mathfrak{a}(\lambda) = 0$ if $|\{ j \in J_{\rhobar}^c: \lambda_j(x_j) \in \{ p-1-x_j, x_j \} \}| < d_\lambda$.
(The converse is also true, by comparing equations~\cite[(81) \& (82)]{BHHMS4}, or alternatively see the proof of \cite[Cor.\,4.4.7]{BHHMS4}.)

\begin{rem}\label{rem:gr-pi2-CM}
Theorem~\ref{thm:gr-pi2} implies that $\gr_\m(\pi_2^\vee)$ is Cohen--Macaulay or zero as $\gr(\Lambda)$-module.
(By \cite[Prop.\,4.4.3]{BHHMS4} and \cite[Cor.\,4.4.5]{BHHMS4}, each nonzero $\mathfrak{a}_1^{i_0}(\lambda)/\mathfrak{a}(\lambda)$ is Cohen--Macaulay, as~the Cohen--Macaulay property is closed under direct summands, shifts in grading, and direct sums.)
\end{rem}

\begin{rem}\label{rem:killed-by-J}
{Theorem~\ref{thm:gr-pi2} shows that $\gr_\m(\pi_2^\vee)$ is killed by the ideal $J$, \ie is an $\ovl R$-module.
This is a priori not obvious.
A similar comment applies to Corollary~\ref{cor:gr-pi'}.}
\end{rem}

\begin{rem}\label{rem:i0-f-1-or-f}
When $i_0 = f$, Theorem~\ref{thm:gr-pi2} is trivially true, because $\pi_2 = 0$ and $\mathfrak{a}_1^{f}(\lambda) = \mathfrak{a}(\lambda)$ for all $\lambda \in \P$.
(By \cite[Lem.\,4.1.4]{BHHMS4} we have
\[
f+1-|J_\lambda| = |J_1|+|J_2|+|J_{\lambda^*}|+1 > |J_1|+|J_2|,
\]
where $\lambda\mto \lambda^*$ is the involution of $\mathscr{P}$ defined in \cite[Def.\,3.62]{BHHMS2}.)
When $i_0 = f-1$, Theorem~\ref{thm:gr-pi2} can be proved as follows, assuming that $\brho$ is only $\max\{9,2f+1\}$-generic.

By \cite[Th.\,4.4.8(ii)]{BHHMS4}, $\pi_2$ is irreducible in this case, so~$\pi_2$ is the principal series $\Ind_{B(K)}^{\GL_2(K)} (\chi_1 \otimes \chi_2\omega^{-1})$ by \cite[Prop.\,10.8]{HuWang2}, where $\ovl\rho \cong \smatr{\chi_1} * 0 {\chi_2}$ and hence
\[
\chi_1|_{I_K} = \omega_f^{\sum_{j=0}^{f-1} (r_j+1) p^j}\quand\chi_2|_{I_K} = 1.
\]
(Here $B(K)$ denotes the Borel subgroup of upper-triangular matrices of $\GL_2(K)$.)

\noindent
\begin{minipage}{\textwidth}\indent
We apply the combinatorial Proposition~\ref{prop:gr-m-semisimple-nonsplit} below (or argue directly) to deduce
\begin{multline}\label{eq:i0-f-1}
\bigoplus_{\lambda \in \P} \chi_\lambda^{-1} \otimes \frac{\mathfrak{a}_1^{i_0}(\lambda)}{\mathfrak{a}(\lambda)}(-d_\lambda)\\[-15pt]
\cong \chi_{\lambda'}^{-1} \!\otimes \!\frac{\ovl R}{(z_j: 0 \le j \le f-1)} \ \oplus \ \chi_{\lambda''}^{-1}\! \otimes \!\frac{\ovl R}{(y_j: 0 \le j \le f-1)},
\end{multline}
\end{minipage}
\par\vspace{6pt}\noindent
where $\lambda',\lambda'' \in \P^\ss$ are given by $\lambda'_j(x_j) = p-3-x_j$, $\lambda''_j(x_j) = x_j+2$ for all $j$. We~calculate $\chi_{\lambda'} = (\chi_2|_{I_K})\omega^{-1} \otimes \chi_1|_{I_K}$ and $\chi_{\lambda''} = \chi_1|_{I_K} \otimes (\chi_2|_{I_K})\omega^{-1}$. We~conclude by \cite[Prop.\,3.76(ii)]{BHHMS2}.
\end{rem}

\begin{lem}\label{lem:lift-graded-homo}

Suppose that $M$ is a graded $\gr(\Lambda)$-module.
Let
\[
N \defeq (\mathfrak{a}_1^{i_0}(\lambda)/\mathfrak{a}(\lambda))(-d_\lambda)
\]
for some $\lambda \in \P$.
Then the natural map
\begin{equation*}
\HOM_{\gr(\Lambda)}(N,M)_0 \to \HOM_{\gr(\Lambda)}(N,M/\gr_{\le-3} M)_0
\end{equation*}
is an isomorphism.
\end{lem}

Recall that $\HOM(N,M)_0$ denotes the graded morphisms $N \to M$ (of~degree 0).

\begin{proof}[Proof. Step 1]
Suppose for the moment that $S$ is a graded ring and $N$ any finitely presented graded $S$-module.
For any subset $D \subset \Z$ we say that \emph{$N$ has relations in degrees $D$} as $S$-module if there exists a graded exact sequence of the form
\[
\bigoplus_{i=1}^n S(-d_i) \to S^{\oplus m} \to N \to 0
\]
with $d_i \in D$ for all $i$ (in~particular, $N$ is generated by its degree 0 part).

We claim that if $N$ has relations in degrees $D$ as $S/I$-module, where $I$ is an ideal of $S$ that is generated by finitely many homogeneous elements $s_i$ whose degrees are contained in $D$, then the same is true as $S$-module.

To see this, by assumption we can find a surjective graded homomorphism
\[
(S/I)^{\oplus m} \to N\to 0,
\]
whose kernel is generated by finitely many homogeneous elements $x_j$ whose degrees are contained in $D$ for all $j$.
Lift each $x_j$ to a homogeneous element $\wt x_j$ of $S^{\oplus m}$ of the same degree. By~composition we have a surjective morphism $S^{\oplus m} \to N \to 0$ of graded $S$-modules.
Its kernel is generated by all $s_i e_k$ (where $e_k$ denotes the standard $\F$-basis of $S^{\oplus m}$) and all $\wt x_j$, as~desired.

\subsubsection*{Step 2}
We show that $N = (\mathfrak{a}_1^{i_0}(\lambda)/\mathfrak{a}(\lambda))(-d_\lambda)$ has relations in degrees $\{-1,-2\}$ as $\gr(\Lambda)$-module.
Since $N$ is a graded $\ovl R/\mathfrak{a}(\lambda)$-module and $\ovl R/\mathfrak{a}(\lambda)$ is obtained from $\gr(\Lambda)$ by quotienting by an ideal generated by homogeneous elements of degrees $-1$ and $-2$, by Step 1 it suffices to show that $N$ has relations in degrees $\{-1,-2\}$ as $\ovl R/\mathfrak{a}(\lambda)$-module.

Note that $j \in J_1 \sqcup J_2$ implies that $t_j = y_jz_j$ and let $d \defeq d_\lambda$ for short. By~interchanging $y_j$ and $z_j$ for some $j$ and permuting $\{0,1,\dots,f-1\}$, we~may assume that
\[
N = (y_{i_1}\cdots y_{i_d}: 0 \le i_1 < \cdots < i_d < e)(-d)
\]
as graded $\ovl R/\mathfrak{a}(\lambda)$-module, for some $e$ such that $0 \le e \le f$.
This module is generated by the elements $X_I \defeq \prod_{i \in I} y_i$ (of~degree 0) for subsets $I \subset E \defeq \{0,1,\dots,e-1\}$ with $|I| = d$. We~claim that the relations are generated by
\begin{equation}\label{eq:relations}
\begin{aligned}
z_i X_I &= 0 && \text{for all $i \in I$}; \\
y_i X_{I'\setminus \{i\}} &= y_j X_{I'\setminus \{j\}} && \text{for all $i \ne j$ in $I' \subset E$, $|I'| = d+1$}.
\end{aligned}
\end{equation}

Note that $\ovl R/\mathfrak{a}(\lambda)$ has as $\F$-basis all monomials of the form $\prod_j w_j^{\ge 0}$ with \hbox{$w_j\!\in\! \{y_j,z_j\}\!\setminus\! \{t_j\}$}.
If $\sum_{I} f_I X_I = 0$ with $f_I \in \ovl R/\mathfrak{a}(\lambda)$, then using relations~\eqref{eq:relations}, without loss of generality, $f_I$ does not contain any $z_i$ ($i \in I$) and $y_i$ ($i \in E \setminus I$, $i < \min I$).
The map $f_I \mto f_I X_I$ is injective for such $f_I$, and moreover for every monomial term in $f_I X_I$, $I$ is the set of $d$ largest elements $i$ of $E$ such that $y_i$ divides~it.
This shows that $f_I = 0$ for all $I$, proving that we have found all relations, and indeed the relations are in degree $-1$.

\subsubsection*{Step 3}
By Step 2 we have an exact sequence\vspace*{-3pt}
\[
\bigoplus_{i=1}^n \gr(\Lambda)(-d_i) \to \gr(\Lambda)^{\oplus m} \to N \to 0
\]
with $d_i \in \{-1,-2\}$ for all $i$. We~get a commutative diagram\vspace*{-3pt}
\begin{multline*}
\xymatrix@R-1pc@C=.5cm{
0 \ar[r] & \HOM_{\gr(\Lambda)}(N,M)_0 \ar[r]\ar[d] & \HOM_{\gr(\Lambda)}(\gr(\Lambda),M)^{\oplus m}_0 \ar[d] \\
0 \ar[r] & \HOM_{\gr(\Lambda)}(N,\ovl M)_0 \ar[r] & \HOM_{\gr(\Lambda)}(\gr(\Lambda),\ovl M)^{\oplus m}_0
}\\
\xymatrix@R-1pc@C=.5cm{
\ar[r]& \bigoplus_{i=1}^n \HOM_{\gr(\Lambda)}(\gr(\Lambda)(-d_i),M)_0 \ar[d]\\
\ar[r] & \bigoplus_{i=1}^n \HOM_{\gr(\Lambda)}(\gr(\Lambda)(-d_i),\ovl M)_0}
\end{multline*}
where $\ovl M \defeq M/\gr_{\le-3} M$. As~$\HOM_{\gr(\Lambda)}(\gr(\Lambda)(-i),M)_0 = M_{i}$ for all $i \in \Z$, the middle and right vertical arrows are isomorphisms, hence so is the left one.
\end{proof}

Fix $n \ge 1$, which we will specify later{, and assume that $\brho$ is $(2n-1)$-generic}.
Recall $\tau \defeq \tau^{(n)} \subset \pi$ from \cite[\S 2.4]{BHHMS4}, so
\[
\tau = \bigoplus_{\lambda \in \P} \tau_\lambda\quad \text{with }\tau_\lambda\defeq \tau_\lambda^{(n)}\text{ and }\soc_I(\tau_\lambda) \cong \chi_\lambda.
\]
Let $\ovl\tau \defeq \tau[\m^{n}]=\pi[\m^{n}]$ (last statement of \cite[Lem.\,2.4.2]{BHHMS4}) and $\ovl\tau_\lambda \defeq \tau_\lambda[\m^{n}]$ for $\lambda \in \P$, so~$\ovl\tau = \bigoplus_{\lambda \in \P} \ovl\tau_\lambda$.

Let $\Theta$ denote the image of $\ovl\tau$ in $\pi_2$. As~$\tau$ is multiplicity-free by \cite[Cor.\,2.4.3(i)]{BHHMS4} (for $r=1$),
we have $\ovl\tau \cap \pi_1 = \bigoplus_{\lambda \in \P} (\ovl\tau_\lambda \cap \pi_1)$ and $\Theta = \bigoplus_{\lambda \in \P} \Theta_\lambda$, where $\Theta_\lambda$ is the image of $\ovl\tau_\lambda$ in $\pi_2$.
For the same reason,\vspace*{-3pt}
\begin{equation}\label{eq:filtrations}
F_{-i} \Theta^\vee = \m^i \ovl\tau^\vee \cap \Theta^\vee = \Bigl(\bigoplus_{\lambda \in \P} \m^i \ovl\tau_\lambda^\vee\Bigr) \cap \Theta^\vee = \bigoplus_{\lambda \in \P} (\m^i \ovl\tau_\lambda^\vee \cap \Theta^\vee) = \bigoplus_{\lambda \in \P} F_{\lambda,-i} \Theta_\lambda^\vee
\end{equation}
for all $i \in \Z_{\ge 0}$, where $F$ (\resp $F_\lambda$) denotes the filtration on $\Theta^\vee$ (\resp $\Theta_\lambda^\vee$) induced from the $\m$-adic filtration on $\ovl\tau^\vee$ (\resp $\ovl\tau_\lambda^\vee$). In~particular, $\gr_F(\Theta^\vee) = \bigoplus_{\lambda \in \P} \gr_{F_\lambda}(\Theta_\lambda^\vee)$ and $F_{-n} \Theta^\vee = 0$.

Suppose that $n \ge 1$ and that $\brho$ is {$(2n-1)$}-generic.
The following lemma determines the submodule structure of $\tau_\lambda^{(n)}[\m^n]$ (and hence of $\pi[\m^n] = \tau^{(n)}[\m^n]$ if $r=1$ by \cite[Cor.\,2.4.3(ii)]{BHHMS4}), since $\bigoplus_{\lambda\in\mathscr{P}}\tau_{\lambda}^{(n)}$ is multiplicity-free by \cite[Cor.\,2.4.3(i)]{BHHMS4} (as $\brho$ is {$(2n-1)$}-generic).

\begin{lem}\label{lem:submodule-structure}
Suppose that $\brho$ is {$(n-1)$-generic}
and keep the above notation.
Suppose that $\lambda \in \P$.
For any $\chi,\chi' \in \JH(\tau^{(n)}_\lambda[\m^n])$ with $\Ext^1_{I/Z_1}(\chi,\chi') \ne 0$, upon perhaps interchanging $\chi$ and $\chi'$, there exist $\ell_j \in \Z$ for $0 \le j \le f-1$, $\ve \in \{\pm1\}$, and $0 \le j_0 \le f-1$ such that
\begin{enumerate}
\item\label{lem:submodule-structurei} $\chi = \chi_\lambda \prod_j \alpha_j^{\ell_j}$ and $\chi' = \chi \alpha_{j_0}^{\ve}$;
\item\label{lem:submodule-structureii} $\ell_j \ge 0$ if $t_j = y_j$, $\ell_j \le 0$ if $t_j = z_j$, and $\sum_j |\ell_j| < n$;
\item\label{lem:submodule-structureiii} $|\ell_{j_0}+\ve| < |\ell_{j_0}|$;
\item\label{lem:submodule-structureiv} $E_{\chi',\chi}$ (the unique nonsplit extension of $\chi$ by $\chi'$, see Section~\ref{sec:preliminaries}) is a subquotient of $\tau^{(n)}_\lambda[\m^n]$.
\end{enumerate}
In particular, either $E_{\chi,\chi'}$ or $E_{\chi',\chi}$ occurs as subquotient of $\tau^{(n)}_\lambda[\m^n]$.
\end{lem}

\begin{proof}
By construction of $\tau^{(n)}_\lambda$ (\cf the proofs of \cite[Lem.\,2.4.1]{BHHMS4} and \cite[Prop.\,9.19]{HuWang2}), as~$\chi\in \tau^{(n)}_\lambda[\m^n]$ we can write $\chi = \chi_\lambda \prod_j \alpha_j^{\ell_j}$ for some $\ell_j \in \Z$ satisfying condition \eqref{lem:submodule-structureii}. As~$\Ext^1_{I/Z_1}(\chi,\chi') \ne 0$ we have $\chi' = \chi \alpha_{j_0}^{\ve}$ for some $\ve \in \{\pm1\}$ and some $j_0$ such that $0 \le j_0 \le f-1$. By~the genericity condition we deduce that condition \eqref{lem:submodule-structureii} holds for $(\ell_0,\dots,\ell_{j_0}+\ve,\dots,\ell_{f-1})$.
(As $\brho$ is {$(n-1)$}-generic, $|\ell_j| < n \le \psfrac{p-1}2$, and
\[
\ve p^{j_0}+\sum_{j=0}^{f-1} \ell_jp^j \equiv \sum_{j=0}^{f-1} \ell'_jp^j \pmod{p^f-1}\quad\text{for integers }|\ell'_j| < \psfrac{p-1}2
\]
implies $(\ell_0,\dots,\ell_{j_0}+\ve,\dots,\ell_{f-1})=(\ell'_0,\dots,\ell'_{f-1})$.)

Interchanging $\chi$ and $\chi'$, if~necessary, we~may assume that \eqref{lem:submodule-structureiii} holds.
Then \eqref{lem:submodule-structureiv} holds, as~the nonsplit extension of $\alpha_{j_0}^{\ell_{j_0}}$ by $\alpha_{j_0}^{\ell_{j_0}+\ve}$ occurs in the $j_0$-th tensor factor defining $\tau^{(n)}_\lambda$.
\end{proof}

\begin{proof}[Proof of Theorem~\ref{thm:gr-pi2}]
By Remark \ref{rem:i0-f-1-or-f} we can assume throughout the proof that $i_0\leq f-2$. Let $N_2'$ denote the right-hand side of the theorem, \ie $N_2' \defeq \bigoplus_{\lambda \in \P} N_{2,\lambda}'$ with
\[
N_{2,\lambda}' \defeq \chi_\lambda^{-1} \otimes \frac{\mathfrak{a}_1^{i_0}(\lambda)}{\mathfrak{a}(\lambda)}(-d_\lambda).
\]

\subsubsection*{Step 1}
We show that $\gr_\m(\pi_2^\vee)/ \ovl\m^3 \onto N_2'/\ovl\m^3$ as graded $\gr(\Lambda)$-modules with compatible $H$-actions.

Consider the commutative diagram
\begin{equation*}
\xymatrix@C=.4cm@R=.7cm{
0 \ar[r] & \bigoplus_{\lambda \in \P} \chi_\lambda^{-1} \otimes \frac{\mathfrak{a}_1^{i_0}(\lambda)}{\mathfrak{a}(\lambda)} \ar[r]\ar^{\cong}[d] & \bigoplus_{\lambda \in \P} \chi_\lambda^{-1} \otimes \frac{\ovl R}{\mathfrak{a}(\lambda)} \ar[r]\ar^{\cong}[d] & \bigoplus_{\lambda \in \P} \chi_\lambda^{-1} \otimes \frac{\ovl R}{\mathfrak{a}_1^{i_0}(\lambda)} \ar[r]\ar^{\cong}[d] & 0 \\
0 \ar[r] & \gr_{F'}(\pi_2^\vee) \ar[r]\ar[d] & \gr(\pi^\vee) \ar[r]\ar@{->>}[d] & \gr(\pi_1^\vee) \ar[r]\ar@{->>}[d] & 0 \\
0 \ar[r] & \gr_F(\Theta^\vee) \ar[r]\ar@{=}[d] & \gr(\ovl\tau^\vee) \ar[r]\ar@{=}[d] & \gr((\ovl\tau\cap \pi_1)^\vee) \ar[r]\ar@{=}[d] & 0 \\
0 \ar[r] & \bigoplus_{\lambda \in \P} \gr_{F_\lambda}(\Theta_\lambda^\vee) \ar[r] & \bigoplus_{\lambda \in \P} \gr(\ovl\tau_\lambda^\vee) \ar[r] & \bigoplus_{\lambda \in \P} \gr((\ovl\tau_\lambda \cap \pi_1)^\vee) \ar[r] & 0
}
\end{equation*}
of graded $\gr(\Lambda)$-modules with compatible $H$-actions, where $F'$ denotes the filtration on $\pi_2^\vee$ induced by the $\m$-adic filtration on $\pi^\vee$ and we recall that $F$ denotes the filtration on $\Theta^\vee$ induced by the $\m$-adic filtration on $\ovl\tau^\vee$.
The top vertical maps are isomorphisms by \cite[Cor.\,4.4.5]{BHHMS4} (see also the proof of \cite[Prop.\,4.4.3]{BHHMS4}).
From $\ovl\tau = \pi[\m^{n}]$ we get $(\ovl\tau\cap \pi_1)[\m^{n}] = \pi_1[\m^{n}]$.
Hence the middle and right vertical maps are isomorphisms in degrees $> -n$,
and so the same is true of the left vertical map.
The middle vertical composition $\bigoplus_{\lambda \in \P} \chi_\lambda^{-1} \otimes (\mathfrak{a}_1^{i_0}(\lambda)/\mathfrak{a}(\lambda)) \to \bigoplus_{\lambda \in \P} \gr(\ovl\tau_\lambda^\vee)$ is an isomorphism in degree 0, respecting the direct sum decomposition (by $H$-equivariance). As~its domain is generated by its degree 0 part as $\gr(\Lambda)$-module, it~follows that the middle vertical composition respects the direct sum decomposition, and hence the same is true for the left and right vertical maps. We~deduce that for each $\lambda \in \P$ the morphism
\begin{equation}\label{eq:N_2-lambda'}
N_{2,\lambda}'(d_\lambda) = \chi_\lambda^{-1} \otimes \frac{\mathfrak{a}_1^{i_0}(\lambda)}{\mathfrak{a}(\lambda)} \to \gr_{F_\lambda}(\Theta_\lambda^\vee)
\end{equation}
of graded $\gr(\Lambda)$-modules is an isomorphism in degrees $> -n$.

We now show that
\begin{equation}\label{eq:compare-fil}
\text{$F_{\lambda,-d_\lambda-i}(\Theta_\lambda^\vee) = \m^i \Theta_\lambda^\vee$ \ \ for any $\lambda \in \P$, $i \ge 0$.}
\end{equation}
To see this, note that if $d_\lambda = 0$, then $\mathfrak{a}_1^{i_0}(\lambda) = \ovl R$, hence the $\lambda$-part of the above commutative diagram shows that the natural map $\gr_{F_\lambda}(\Theta_\lambda^\vee) \hto \gr(\ovl\tau_\lambda^\vee)$ is an isomorphism, which implies that the natural map $\Theta_\lambda^\vee \hto \ovl\tau_\lambda^\vee$ of filtered $\Lambda$-modules is an isomorphism, as~desired.
Suppose now that $d_\lambda > 0$. We~first obtain from the previous diagram the following diagram:
\begin{equation*}
\xymatrix{
0 \ar[r] & \chi_\lambda^{-1} \otimes \frac{\mathfrak{a}_1^{i_0}(\lambda) + \m^n}{\mathfrak{a}(\lambda)+\m^n} \ar[r]\ar^{\cong}[d] & \chi_\lambda^{-1} \otimes \frac{\ovl R}{\mathfrak{a}(\lambda)+\m^n} \ar[r]\ar^{\cong}[d] & \chi_\lambda^{-1} \otimes \frac{\ovl R}{\mathfrak{a}_1^{i_0}(\lambda)+\m^n} \ar[r]\ar^{\cong}[d] & 0 \\
0 \ar[r] & \gr_{F_\lambda}(\Theta_\lambda^\vee) \ar[r] & \gr(\ovl\tau_\lambda^\vee) \ar[r] & \gr((\ovl\tau_\lambda \cap \pi_1)^\vee) \ar[r] & 0
}
\end{equation*}
By the definition of $\mathfrak{a}(\lambda)$ the middle isomorphism shows that\vspace*{-3pt}
\begin{equation*}\textstyle
\JH(\ovl\tau_\lambda^\vee) = \big\{ \chi_\lambda^{-1} \prod_j \alpha_j^{\ell_j}: \ell \in L \big\},
\end{equation*}
where\vspace*{-3pt}
\begin{equation*}\textstyle
L \defeq \big\{ \ell = (\ell_j)_{j=0}^{f-1}: \text{$\ell_j \le 0$ if $t_j = y_j$, $\ell_j \ge 0$ if $t_j = z_j$},\ \sum_j |\ell_j| < n \big\},
\end{equation*}
as well as\vspace*{-3pt}
\begin{equation*}\textstyle
\JH(\m^i \ovl\tau_\lambda^\vee) = \big\{ \chi_\lambda^{-1} \prod_j \alpha_j^{\ell_j}: \ell \in L,\ i \le \sum_j |\ell_j| \big\}.
\end{equation*}
By the definition of $\mathfrak{a}_1^{i_0}(\lambda)$ the left isomorphism shows that
\begin{equation*}\textstyle
\JH(\Theta_\lambda^\vee) = \big\{ \chi_\lambda^{-1} \prod_j \alpha_j^{\ell_j}: \ell \in L,\; |\{j \in J_1: \ell_j > 0\}|+|\{j \in J_2: \ell_j < 0\}| \ge d_\lambda \big\},
\end{equation*}
where we recall that\vspace*{-3pt}
\[
J_1 = \{ j \in J_{\rhobar}^c: \lambda_j(x_j) = p-1-x_j \}\quand J_2 = \{ j \in J_{\rhobar}^c: \lambda_j(x_j) = x_j \}.
\]
In particular, $\Theta_\lambda^\vee \subset \m^{d_\lambda} \ovl\tau_\lambda^\vee$ and hence $\m^i \Theta_\lambda^\vee \subset \Theta_\lambda^\vee \cap \m^{d_\lambda+i} \ovl\tau_\lambda^\vee$ for all $i \ge 0$.
Conversely, to show $\Theta_\lambda^\vee \cap \m^{d_\lambda+i} \ovl\tau_\lambda^\vee \subset \m^i \Theta_\lambda^\vee$, by multiplicity-freeness it suffices to show that $\JH(\Theta_\lambda^\vee) \cap \JH(\m^{d_\lambda+i} \ovl\tau_\lambda^\vee) \subset \JH(\m^i \Theta_\lambda^\vee)$.
Take $\chi \defeq \chi_\lambda^{-1} \prod_j \alpha_j^{\ell_j}$ with\vspace*{-3pt}
\[
\ell \in L,\quad |\{j \in J_1: \ell_j > 0\}|+|\{j \in J_2: \ell_j < 0\}| \ge d_\lambda,\quand d_\lambda+i \le \sum_j |\ell_j|.
\]
With the help of Lemma \ref{lem:submodule-structure} it is easy to show that there exist characters $\chi_{i'} \in \JH(\Theta_\lambda^\vee)$ ($0 \le i' \le i$) with $\chi_0 \defeq \chi$ and such that the unique nonsplit extension $E_{\chi_{i'-1},\chi_{i'}}$ (of~$\chi_{i'}$ by $\chi_{i'-1}$) occurs as subquotient of $\Theta_\lambda^\vee$ for all $i'$ such that $0 < i' \le i$.
(If $i > 0$ then we find $\chi_1$ as follows: if there is $j \notin J_1 \sqcup J_2$ such that $\ell_j \ne 0$, choose such a $j$;
otherwise, choose $j$ such that $|\ell_j| > 0$ and is as small as possible.
Then $\chi_1 \defeq \chi_0 \alpha_j^{-\sgn(\ell_j)}$ is still an element of $\JH(\Theta_\lambda^\vee)$, and we have decreased $\sum_j |\ell_j|$ by~$1$.
Proceed inductively to find all $\chi_{i'}$.) We deduce that $\chi$ occurs in $\rad^i \Theta_\lambda^\vee = \m^i \Theta_\lambda^\vee$, proving~\eqref{eq:compare-fil}.

We now let $n \defeq i_0+4$. As~$d_\lambda \le i_0+1 < n-2$ we obtain from~\eqref{eq:N_2-lambda'} and~\eqref{eq:compare-fil} an isomorphism of graded $\gr(\Lambda)$-modules
\begin{equation*}
N_{2,\lambda}'/\ovl\m^3 = N_{2,\lambda}'/\gr_{\le -3} N_{2,\lambda}' \cong \gr_{F_\lambda}(\Theta_\lambda^\vee(-d_\lambda))/\gr_{F_\lambda,\le -3}(\Theta_\lambda^\vee(-d_\lambda)) = \gr(\Theta_\lambda^\vee)/\ovl\m^3.
\end{equation*}
Hence\vspace*{-3pt}
\begin{equation}
\label{eq:grad2:surj}
\gr(\pi_2^\vee)/\ovl\m^3 \onto \gr(\Theta^\vee)/\ovl\m^3 \cong \bigoplus_{\lambda \in \P} \gr(\Theta_\lambda^\vee)/\ovl\m^3 \cong N_2'/\ovl\m^3,
\end{equation}
as desired.

\subsubsection*{Step 2} We show that $\gr_\m(\pi_2^\vee)/ \ovl\m^3 \cong N_2'/\ovl\m^3$ as graded $\gr(\Lambda)$-modules with compatible $H$-actions.
From the cohomology long exact sequence we get
\begin{multline}\label{eq:coho-seq}
0 \to \coker\bigl(\Tor_1^\Lambda(\Lambda/\m^3,\pi^\vee) \to \Tor_1^\Lambda(\Lambda/\m^3,\pi_1^\vee)\bigr) \to \pi_2^\vee/\m^3\\
\to \pi^\vee/\m^3 \to \pi_1^\vee/\m^3 \to 0.
\end{multline}
We let
\begin{equation*}
C \defeq \coker\bigl(\Tor_1^\Lambda(\Lambda/\m^3,\pi^\vee) \to \Tor_1^\Lambda(\Lambda/\m^3,\pi_1^\vee)\bigr)
\end{equation*}
and give it the induced filtration as quotient of $\Tor_1^\Lambda(\Lambda/\m^3,\pi_1^\vee)$.

First we show that $\gr(C)$ is a subquotient of
\begin{equation}\label{eq:coker-Tor1-gr}
C' \defeq \coker\bigl(\Tor_1^{\gr(\Lambda)}(\gr(\Lambda)/\ovl\m^3,\gr_{\m}(\pi^{\vee})) \to \Tor_1^{\gr(\Lambda)}(\gr(\Lambda)/\ovl\m^3,\gr_{\m}(\pi_1^{\vee}))\bigr).
\end{equation}
Notice that $\gr(C)$ is a quotient of
\begin{equation*}
\coker\bigl(\gr(\Tor_1^{\Lambda}(\Lambda/\m^3,\pi^{\vee})) \to \gr(\Tor_1^{\Lambda}(\Lambda/\m^3,\pi_1^{\vee}))\bigr),
\end{equation*}
because if we have a filtered exact sequence $X \to Y \to C \to 0$ with $C$ carrying the induced filtration, then $\coker(\gr(X) \to \gr(Y))$ surjects onto $\gr(C)$ by \cite[Th.\,I.4.2.4(1)]{LiOy}.
Then, as~in the proof of \cite[Prop.\,2.4.9]{BHHMS4}, we~consider the morphism of spectral sequences that converges to this morphism:
\begin{equation*}
\begin{gathered}
\xymatrix{E_i^r\ar@{=>}[r]\ar[d]&\Tor_i^{\Lambda}(\Lambda/\m^3,\pi^{\vee})\ar[d]\\
E'^r_i\ar@{=>}[r]&\Tor_i^{\Lambda}(\Lambda/\m^3,\pi_1^{\vee}).}
\end{gathered}
\end{equation*}
(Referring to that proof, we~have $E_i^0 = \gr(\Lambda/\m^3 \otimes_\Lambda M_i) \cong \gr(\Lambda/\m^3) \otimes_{\gr(\Lambda)} \gr(M_i)$ by \cite[Lem.\,I.6.14]{LiOy}, so~$E_i^1 = \Tor_i^{\gr(\Lambda)}(\gr(\Lambda)/\ovl\m^3,\gr_{\m}(\pi^{\vee}))$.)
Assumption~\ref{subsec:assumptions}\eqref{it:assum-v} says that $\dim_\F E_1^\infty = \dim_\F E_1^1$. It~easily follows that $\coker(E_1^{r+1} \to E'^{r+1}_1)$ is a subquotient of $\coker(E_1^r \to E'^r_1)$ for any $r \ge 1$ (recall that $E_1'^{r+1}$ is a subquotient of $E_1'^r$, while $E^{r+1}_1=E^{r}_1$ by the preceding sentence).
This implies the claim by taking $r$ sufficiently large.

From the sequence~\eqref{eq:coho-seq} we see that
\begin{equation}\label{eq:dim-C}
\begin{aligned}
\dim_\F(C) &= \dim_\F(\pi_2^\vee/\m^3)-\dim_\F(\pi^\vee/\m^3)+\dim_\F(\pi_1^\vee/\m^3) \\
& = \dim_\F(\gr_{\m}(\pi_2^\vee)/\ovl\m^3)-\dim_\F(\gr(\pi^\vee)/\ovl\m^3)+\dim_\F(\gr(\pi_1^\vee)/\ovl\m^3).
\end{aligned}
\end{equation}
By Step 1 we know that
\begin{equation}\label{eq:dim-gr-pi2}
\begin{aligned}
\dim_\F(\gr_{\m}(\pi_2^\vee)/\ovl\m^3) &\ge \dim_\F(N_2'/\ovl\m^3) = \sum_{\lambda \in \P} \dim_\F(N_{2,\lambda}'/\ovl\m^3)\\
&= \sum_{\lambda \in \P} \dim_\F(N_{2,\lambda}/\ovl\m^3) = \dim_\F(\gr_{F'}(\pi_2^\vee)/\ovl\m^3),
\end{aligned}
\end{equation}
where $N_{2,\lambda} \defeq \chi_\lambda^{-1} \otimes (\mathfrak{a}_1^{i_0}(\lambda)/\mathfrak{a}(\lambda)) = N_{2,\lambda}'(d_\lambda)$ and we used \cite[Cor.\,4.4.5]{BHHMS4} for the last equality.
Combining equations~\eqref{eq:dim-C}, \eqref{eq:dim-gr-pi2} together with the fact that $\gr(C)$ is a subquotient of $C'$ (\cf \eqref{eq:coker-Tor1-gr}) we obtain
\begin{equation}
\begin{aligned}\label{eq:ineq}
\dim_\F(C') &\ge \dim_\F(C)\\
&\ge \dim_\F(\gr_{F'}(\pi_2^\vee)/\ovl\m^3)
-\dim_\F(\gr(\pi^\vee)/\ovl\m^3)+\dim_\F(\gr(\pi_1^\vee)/\ovl\m^3).
\end{aligned}
\end{equation}
The exact sequence
\begin{equation*}
0 \to C' \to \gr_{F'}(\pi_2^\vee)/\ovl\m^3 \to \gr(\pi^\vee)/\ovl\m^3 \to \gr(\pi_1^\vee)/\ovl\m^3 \to 0
\end{equation*}
shows that equality holds in~\eqref{eq:ineq}, and hence in~\eqref{eq:dim-gr-pi2}. As~equality holds in~\eqref{eq:dim-gr-pi2}, the surjection $\gr_\m(\pi_2^\vee)/ \ovl\m^3 \onto N_2'/\ovl\m^3$ in \eqref{eq:grad2:surj} has to be an isomorphism.

\subsubsection*{Step 3} By Lemma~\ref{lem:lift-graded-homo} we get a graded morphism (of~degree $0$) $N_2' \to \gr_\m(\pi_2^\vee)$, which has to be surjective by the graded Nakayama lemma.
Recall that $N_2'$ is Cohen--Macaulay by Remark~\ref{rem:gr-pi2-CM}.

By Step 4 of the proof of \cite[Prop.\,4.4.3]{BHHMS4} we have $\mathcal{Z}(N_2') = \mathcal{Z}(N_2^{i_0}) = \mathcal{Z}(\gr_\m(\pi_2^\vee))$.
Using the same argument as in the last paragraph of the proof of \cite[Prop.\,4.4.3]{BHHMS4} (\ie $N_2'$ is Cohen--Macaulay and the two modules have the same cycle), we~deduce that the morphism $N_2' \onto \gr_\m(\pi_2^\vee)$ is an isomorphism.
\end{proof}
\setcounter{prop}{9}
For any $\lambda' \in \P^\ss$ we let $\fa^\ss(\lambda')$ denote the ideal of $R$ generated by all $t_j^\ss = t_j^\ss(\lambda') \in \{y_j,z_j,y_jz_j\}$, which are defined as in~\eqref{eq:id:al} but for the Galois representation $\brho^\ss$. Recall $\fa^\ss(\lambda') \supset \ker(R \to \ovl R)$, so~we often think of it as ideal of $\ovl R$.

\begin{prop}\label{prop:gr-m-semisimple-nonsplit}
For any $-1 \le i_0 \le f-1$ we have an isomorphism\vspace*{-3pt}
\begin{equation*}
\bigoplus_{\lambda\in\P}\chi_{\lambda}^{-1}\otimes \frac{\fa_1^{i_0}(\lambda)}{\fa_1^{i_0+1}(\lambda)}(-d_\lambda) \cong \bigoplus_{\lambda'\in\P^\ss,\; \ell(\lambda') = i_0+1}\chi_{\lambda'}^{-1}\otimes \frac{\ovl R}{\fa^\ss(\lambda')}
\end{equation*}
of graded $\gr(\Lambda)$-modules with compatible $H$-actions.
\end{prop}
\begin{proof}
Fix $\lambda\in\P$. As~usual, set
\[
J_1 \defeq \{ j \in J_{\rhobar}^c: \lambda_j(x_j) = p-1-x_j \}\quand J_2 \defeq \{ j \in J_{\rhobar}^c: \lambda_j(x_j) = x_j \},
\]
and let $J \defeq J_1 \sqcup J_2$. We~show that
\begin{equation}\label{eq:lambda-term}
\chi_{\lambda}^{-1}\otimes \frac{\fa_1^{i_0}(\lambda)}{\fa_1^{i_0+1}(\lambda)}(-d_\lambda) \cong \bigoplus_{\substack{J' \subset J\\ |J_\lambda|+|J'|=i_0+1}}\chi_{\lambda'(J')}^{-1}\otimes \frac{\ovl R}{\fa^\ss(\lambda'(J'))},
\end{equation}
where $\lambda'(J') \in \P^\ss$ is defined by
\begin{equation*}
\lambda'(J')_j(x_j) \defeq
\begin{cases}
p-3-x_j & \text{if $j \in J_1' \defeq J' \cap J_1$,} \\
x_j+2 & \text{if $j \in J_2' \defeq J' \cap J_2$,} \\
\lambda_j(x_j) & \text{otherwise.}
\end{cases}
\end{equation*}

It is easy to check that this implies the proposition, by taking a direct sum over all $\lambda \in \P$.

If $i_0+1-|J_\lambda| < 0$, then~\eqref{eq:lambda-term} trivially holds: by definition, $\fa_1^{i_0}(\lambda) = \fa_1^{i_0+1}(\lambda) = \ovl R$, so~both sides are zero.

Suppose that $i_0+1-|J_\lambda| \ge 0$, so~$d_\lambda = i_0+1-|J_\lambda|$.
Note that
\[
\chi_{\lambda'(J')}^{-1} = \chi_\lambda^{-1} \prod_{j \in J_1'}\alpha_j \prod_{j \in J_2'}\alpha_j^{-1},
\]
and $\prod_{j \in J_1'}\alpha_j \prod_{j \in J_2'}\alpha_j^{-1}$ gives the action of $H$ on the degree $d_\lambda$ polynomial
\[
p_{J'} \defeq \prod_{j \in J_1'}y_j \prod_{j \in J_2'}z_j \in R \cong \F[y_j,z_j: 0 \le j \le f-1].
\]
Therefore, by twisting both sides by $\chi_\lambda(d_\lambda)$, it~suffices to show that
\begin{equation*}
\frac{\fa_1^{i_0}(\lambda)}{\fa_1^{i_0+1}(\lambda)} \cong \bigoplus_{\substack{J' \subset J\\ |J'|=d_\lambda}} p_{J'}\otimes \frac{\ovl R}{\fa^\ss(\lambda'(J'))}.
\end{equation*}
We have
\begin{equation*}
\frac{\fa_1^{i_0}(\lambda)}{\fa_1^{i_0+1}(\lambda)} \cong \frac{I(J_1,J_2,d_\lambda)+\fa(\lambda)}{I(J_1,J_2,d_\lambda+1)+\fa(\lambda)} \cong \frac{I(J_1,J_2,d_\lambda)}{I(J_1,J_2,d_\lambda+1)+I(J_1,J_2,d_\lambda) \cap \fa(\lambda)}
\end{equation*}
(recall the ideals $I(J_1,J_2,d_\lambda)$ from \cite[Def.\,4.2.4]{BHHMS4}).
Note that $I(J_1,J_2,d_\lambda) = (p_{J'}: J' \subset J,\; |J'|=d_\lambda)$ and $\fa(\lambda) = (t_j: 0 \le j \le f-1)$ with $t_j = y_jz_j$ for all $j \in J$ (see equation \eqref{eq:id:al}). By~\cite[Prop.\,1.2.1]{herzog-hibi} the ideal $I(J_1,J_2,d_\lambda) \cap \fa(\lambda)$ is generated by the monomials $p_{J'} z_j$ ($j \in J_1'$), $p_{J'} y_j$ ($j \in J_2'$), and $p_{J'} t_j$ ($j \notin J'$), where $J' \subset J$ runs through all subsets with $|J'|=d_\lambda$.
Hence the ideal $I(J_1,J_2,d_\lambda+1)+I(J_1,J_2,d_\lambda) \cap \fa(\lambda)$ is generated by the monomials
\[
p_{J'} z_j\; (j \in J_1' \sqcup (J_2\setminus J_2')), \quad p_{J'} y_j\; (j \in J_2' \sqcup (J_1\setminus J_1')), \quand p_{J'} t_j\; (j \notin J),
\]
where again $J' \subset J$ runs through all subsets with $|J'|=d_\lambda$.
Since, from equation \eqref{eq:id:al} and the definition of $\lambda'(J')$,
\[
\fa^\ss(\lambda'(J')) = \fa(\lambda) + (z_j: j \in J_1' \sqcup (J_2\setminus J_2')) + (y_j: j \in J_2' \sqcup (J_1\setminus J_1')),
\]
we deduce that
\begin{equation*}
I(J_1,J_2,d_\lambda+1)+I(J_1,J_2,d_\lambda) \cap \fa(\lambda) = \sum_{J' \subset J,\; |J'|=d_\lambda} p_{J'} \cdot \fa^\ss(\lambda'(J')).
\end{equation*}
In particular,
\[
\frac{I(J_1,J_2,d_\lambda)}{I(J_1,J_2,d_\lambda+1)+I(J_1,J_2,d_\lambda) \cap \fa(\lambda)} \cong \frac{\sum_{J' \subset J,\;
|J'|=d_\lambda} p_{J'}R}{\sum_{J' \subset J,\; |J'|=d_\lambda} p_{J'} \fa^\ss(\lambda'(J'))},
\]
so for each index $J'$, multiplication induces a homomorphism
\begin{equation*}
p_{J'}\otimes \frac{R}{\fa^\ss(\lambda'(J'))} \to \frac{I(J_1,J_2,d_\lambda)}{I(J_1,J_2,d_\lambda+1)+I(J_1,J_2,d_\lambda) \cap \fa(\lambda)}
\end{equation*}
and passing to the direct sum induces a surjective homomorphism
\begin{equation*}
\bigoplus_{\substack{J' \subset J\\ |J'|=d_\lambda}} p_{J'}\otimes \frac{R}{\fa^\ss(\lambda'(J'))} \onto \frac{I(J_1,J_2,d_\lambda)}{I(J_1,J_2,d_\lambda+1)+I(J_1,J_2,d_\lambda) \cap \fa(\lambda)},
\end{equation*}
which we need to show is an isomorphism.
Suppose that we have an element $f = (p_{J'} \otimes [f_{J'}])_{J'}$ in the kernel, or equivalently that $\sum_{J'} p_{J'}f_{J'} = \sum_{J'} p_{J'}g_{J'}$ in $R$ for some $g_{J'} \in \fa^\ss(\lambda'(J'))$. By~replacing $f_{J'}$ by $f_{J'}-g_{J'}$ we may assume that $g_{J'} = 0$ for all $J'$, \ie that $\sum_{J'} p_{J'}f_{J'} = 0$.
Fix $J'$ and let
\[
\fb(J') \defeq (z_j: j \in J_2\setminus J_2') + (y_j: j \in J_1\setminus J_1') \subset \fa^\ss(\lambda'(J')).
\]
From $p_{J'}f_{J'} = - \sum_{J'' \ne J'} p_{J''}f_{J''}$ we deduce that $p_{J'}f_{J'} \in \fb(J')$. As~multiplication by $p_{J'}$ is injective on $R/\fb(J')$ (a polynomial ring), it~follows that
\[
f_{J'} \in \fb(J') \subset \fa^\ss(\lambda'(J')),
\]
so~$f = 0$, as~desired.
\end{proof}
\begin{thebibliography}{99}
\bibitem[BDJ10]{BDJ} K. Buzzard, F. Diamond \& F. Jarvis - “On Serre’s conjecture for mod $\ell $ Galois representations over totally real fields”, Duke Math. J. 155 (2010) no. 1, p. 105-161.
\bibitem[BHH+23]{BHHMS1} C. Breuil, F. Herzig, Y. Hu, S. Morra \& B. Schraen - “Gelfand-Kirillov dimension and mod $p$ cohomology for $\rm GL_2$”, Invent. Math. 234 (2023) no. 1, p. 1-128.
\bibitem[BHH+25a]{BHHMS2} C. Breuil, F. Herzig, Y. Hu, S. Morra \& B. Schraen - Conjectures and results on modular representations of ${\rm GL}_n({K})$ for a $p$-adic field ${K}$, Mem. Amer. Math. Soc., vol. 315, no. 1598, American Mathematical Society, Providence, RI, 2025.
\bibitem[BHH+25b]{BHHMS4} C. Breuil, F. Herzig, Y. Hu, S. Morra \& B. Schraen - “Finite length for unramified $\mathrm{GL}_2$”, 2025, arXiv:2501.03644.
\bibitem[BP12]{BP} C. Breuil \& V. Paškūnas - Towards a modulo $p$ Langlands correspondence for ${\rm GL}_2$, Mem. Amer. Math. Soc., vol. 216, no. 1016, American Mathematical Society, Providence, RI, 2012.
\bibitem[Bre14]{breuil-buzzati} C. Breuil - “Sur un problème de compatibilité local-global modulo $p$ pour ${\rm GL}_2$”, J. reine angew. Math. 692 (2014), p. 1-76.
\bibitem[Eme10]{emerton-ordII} M. Emerton - “Ordinary parts of admissible representations of $p$-adic reductive groups II. Derived functors”, in Représentations $p$-adiques de groupes $p$-adiques III : méthodes globales et géométriques, Astérisque, vol. 331, Société Mathématique de France, Paris, 2010, p. 403-459.
\bibitem[HH11]{herzog-hibi} J. Herzog \& T. Hibi - Monomial ideals, Graduate Texts in Math., vol. 260, Springer-Verlag London, Ltd., London, 2011.
\bibitem[HW22]{HuWang2} Y. Hu \& H. Wang - “On the mod $p$ cohomology for $\rm GL_2$: the non-semisimple case”, Camb. J. Math. 10 (2022) no. 2, p. 261-431.
\bibitem[Hu10]{yongquan-algebra} Y. Hu - “Sur quelques représentations supersingulières de ${\rm GL}_2(\mathbb{Q}_{p^f})$”, J. Algebra 324 (2010) no. 7, p. 1577-1615.
\bibitem[Koh17]{Ko} J. Kohlhaase - “Smooth duality in natural characteristic”, Adv. Math. 317 (2017), p. 1-49.
\bibitem[LvO96]{LiOy} H. Li \& F. van Oystaeyen - Zariskian filtrations, Kluwer Monographs in Mathematics, vol. 2, Kluwer Academic Publishers, Dordrecht, 1996.
\end{thebibliography}
\end{document}
